\documentclass[12pt, a4paper]{article}
\usepackage{graphicx} % Required for inserting images
\graphicspath{{images/}{images-desafios/}}
\usepackage[a4paper, total={7.5in, 10in}]{geometry}
\usepackage{float} % pra conseguir usar o \begin{figure}[H]
\usepackage{indentfirst} % pra fazer tab n 1 paragrafos dos capitulos
% package pra tabela
\usepackage{multirow}
\usepackage{array}

%pacote pra codigo
\usepackage{listings} 
\usepackage{xcolor}
\definecolor{codegreen}{rgb}{0,0.6,0}
\definecolor{codegray}{rgb}{0.5,0.5,0.5}
\definecolor{codepurple}{rgb}{0.58,0,0.82}
\definecolor{backcolour}{rgb}{0.95,0.95,0.92}

\lstdefinestyle{myCodeStyle}{
	backgroundcolor=\color{backcolour},   
	commentstyle=\color{codegreen},
	keywordstyle=\color{cyan},
	numberstyle=\tiny\color{codegray},
	stringstyle=\color{codepurple},
	basicstyle=\ttfamily\footnotesize,
	breakatwhitespace=false,         
	breaklines=true,                 
	captionpos=b,                    
	keepspaces=true,                 
	%numbers=left,                    
	numbersep=5pt,                  
	showspaces=false,                
	showstringspaces=false,
	showtabs=false,                  
	tabsize=2
}
\lstset{style=myCodeStyle}
\renewcommand{\lstlistingname}{Código}%Muda o termo "Listing":  Listing -> Código
\renewcommand{\lstlistlistingname}{Lista de \lstlistingname s}% List of Listings -> Lista de Códigos


\newcommand{\unit}[1]{\ensuremath{\, \mathrm{#1}}} % comando pra tirar do math mode a unidade de medida

\def\tamanhoGrafico{0.8}

\newcommand{\ohm}{\unit{\Omega} }
\newcommand{\kohm}{\unit{k\Omega}}

\title{
	Laboratório Digital A - PCS3335
	\\ Planejamento da Experiência 3
}

\author{
	Alunos:
	\\ Rafael Hipólito Rodrigues - N.USP: 14604308
	\\ {{PERSON-2}} - N.USP: {{PERSON-ID-2}}
	\\ \\
	Turma 3 - Bancada A6
	\\ Professor: Bruno de Carvalho Albertini
}

%\date{}

\begin{document}
	
	\maketitle
	
	\section*{Planejamento}
	
	Elaboração de uma entidade semaforo que alterar entre verde, amarelo e vermelho utilizando uma Maquina de Estados Finitos com clock de 1.8432MHz.
	\par A figura \ref{fig:fsmExp3} mostra o Diagrama da Maquina de Estados Finitos da entidade semaforo.
	
	\begin{figure}[H]
		\centering
		\includegraphics[width=0.7\linewidth]{images/fsmExp3.pdf}
		\caption{Diagrama da Maquina de Estados Finitos do semaforo.}
		\label{fig:fsmExp3}
	\end{figure}
	
	\begin{itemize}
		\item verde, amarelo e vermelho são as saídas da entidade.
		\item reset é a entrada reset da entidade.
		\item sig\_verde, sig\_amarelo e sig\_vermelho são os sinais que são ativados quando o contador de seu
		respectivo estado atinge o valor desejado na contagem dos segundos.
		\item sig\_reset é o sinal que é ativado quando o estado atual é \texttt{Sreset}.
	\end{itemize}

	
	\section*{Diagramas RTL}
	
	\begin{figure}[H]
		\centering
		\includegraphics[width=1\linewidth]{images/semaforoRafRTL.pdf}
		\caption{Diagrama RTL do semaforo.}
		\label{fig:semaforoRafRTL}
	\end{figure}
	
	\section*{Simulação}
	
	Simulação da entidade usando GHDL e GTKWave:
	
	\begin{figure}[H]
		\centering
		\includegraphics[width=1\linewidth]{images/vcd6gb.png}
		\caption{Resultado da simalação no GTKWave.}
		\label{fig:vcd6gb}
	\end{figure}
	
	\section*{Pinagem}
	
	\begin{table}[H]
		\begin{center}
			\begin{tabular}{|l|l|}
				\hline
				amarelo  & PIN\_AA22 \\ \hline
				clk      & ip\_pll    \\ \hline
				reset    & PIN\_U7   \\ \hline
				verde    & PIN\_U21  \\ \hline
				vermelho & PIN\_W21  \\ \hline
			\end{tabular}
		\end{center}
		
	\end{table}
	
	\section*{Testes}
	
	\begin{itemize}
		\item começa sinal verde, apos 5s, começa sinal amarelo.
	\end{itemize}

	
	
\end{document}
